\section{ZFS}\label{zfs}

In ZFS si segue questo principio base: \emph{Un blocco in uso non viene
mai sovrascritto.}\footnote{\label{fn:ref_29}OpenZFS Documentation Project, \emph{OpenZFS
  System Administration and Architecture},
  https://openzfs.github.io/openzfs-docs/}

Modificare dati vuol dire dunque scrivere il contenuto in una nuova
porzione di memoria, ed il filesystem dovrà essere modificato per
puntare a questa nuova area.

Questo principio, permette di implementare e gestire facilmente feature
avanzate come snapshots ed il self-healing dei dati.

In ZFS lo storage è organizzato all'interno di un pool (denominato
\emph{zpool}), strutturato come un albero di blocchi:

\begin{itemize}
\tightlist
\item
  le foglie contengono i blocchi di dati effettivi;
\item
  i nodi intermedi sono puntatori indiretti a blocchi (\textbf{block
  pointers});
\item
  l'\textbf{uberblock} è il nodo radice dell'albero, da cui è navigabile
  l'intero filesystem.
\end{itemize}

Un block pointer tiene traccia di un insieme di blocchi memorizzandone
dimensione e \textbf{checksum}. Ciascun nodo intermedio valida
crittograficamente i propri figli. Di conseguenza, modificare dati nel
filesystem comporta il ricalcolo delle checksum e la riscrittura della
catena dei nodi padri a ritroso fino alla radice, allocando nuovi
blocchi per ogni livello.

Al termine della catena di modifiche, viene generato un nuovo
\(\text{uberblock'}\), che fa riferimento al nuovo albero con i dati
aggiornati e i puntatori consistenti. L'intera transizione avviene in
modo \textbf{atomico}.

Il costo computazionale per una modifica è proporzionale all'altezza
dell'albero, non alla dimensione assoluta dello storage pool. Questa
proprietà permette di mantenere contemporaneamente valido sia il nuovo
\(\text{uberblock'}\) sia il precedente \(\text{uberblock}\).

ZFS segue un modello transazionale: le operazioni di scrittura vengono
aggregate in gruppi di transazione (\emph{transaction groups},
\passthrough{\lstinline!txg!}) applicati atomicamente su disco.

\subsection{Vantaggi di Pool e
transazioni}\label{vantaggi-di-pool-e-transazioni}

\emph{Gli snapshot del filesystem vengono eseguiti con complessità
temporale \(O(1)\)}: creare uno snapshot consiste semplicemente
nell'impedire che i blocchi foglia dell'albero al momento dello snapshot
vengano deallocati o sovrascritti. Per preservare l'intero stato del
filesystem è sufficiente conservare il riferimento
all'\(\text{uberblock}\) congelato.

\begin{quote}
Uno snapshot in ZFS è una vista read-only del dataset congelata a un
istante preciso. La sua creazione è un'operazione atomica, istantanea e
a costo zero di spazio disco iniziale.
\end{quote}

\textbf{Integrità dei dati e self-healing}: ZFS rileva corruzioni
silenti confrontando le checksum memorizzate nei puntatori padri con il
contenuto effettivo dei blocchi. Qualora il pool sia configurato con
ridondanza (mirroring o RAID-Z), il filesystem è in grado di correggere
automaticamente il dato corrotto attingendo dalla copia integra.

\subsection{Considerazioni di Efficienza e
Frammentazione}\label{considerazioni-di-efficienza-e-frammentazione}

La natura Copy-on-Write implica che l'aggiornamento di dati preesistenti
allochi blocchi in posizioni non contigue dello storage pool,
frammentando progressivamente la sequenzialità fisica dei dati su disco
(\emph{fragmentation}). L'efficienza dell'allocatore di ZFS dipende
pertanto dalla percentuale di spazio libero residuo nel pool e dal
livello di frammentazione complessivo.

Inoltre, qualsiasi operazione di scrittura (inclusa la cancellazione di
file) richiede l'allocazione di nuovi metadati. Nel caso di scritture
parziali e casuali di piccole dimensioni, il costo di riallocazione dei
metadati e il ricalcolo delle checksum possono risultare significativi.
